\section{Polytope faces and exact kissing configurations}
\label{sec:polygeom}

The survey names the Birkhoff polytope and the permutohedron, asks how
many faces of each dimension a polytope can have, and discusses kissing
numbers. For these questions, rational points and finite vertex lists
already support exact theorems. A rational convex combination is a list
of nonnegative integer weights \(w_i\) with common denominator
\(D=\sum_i w_i>0\). For a finite rational vertex list \(V\), a
nonempty exposed face is the convex hull of those vertices maximizing
a supplied rational linear functional; the zero functional includes
the whole polytope. Maximizers over the convex hull have precisely this
description: a weighted average of numbers \(\le M\) equals \(M\)
exactly when every positive-weight term equals \(M\).
Dimension means the rank over \(\mathbb Q\) of the differences from one
vertex of the face. In the proofs below explicit bases compute these ranks.

\subsection{Assignment matrices, faces, and optimization}

\begin{theorem}[Rational Birkhoff--von Neumann theorem]\label{result54}
Let \(A\) be a nonnegative rational \(n\times n\) matrix, \(n\ge1\),
with all row and column sums equal to \(1\). There is an explicit convex
decomposition into permutation matrices with at most \(n^2\) terms.
With common denominator \(D>0\), the terms can all have weights with
denominator \(D\). The \(n!\) permutation matrices are precisely the
extreme points. Maximizing a rational linear objective on this polytope
is equivalent to maximizing it over permutations.
\end{theorem}
\begin{proof}
Clear denominators, obtaining a nonnegative integer matrix with common
row and column sum \(s=D\). Whenever \(s>0\), join row \(i\) to column
\(j\) if the entry is positive. For any subset \(U\) of rows, all of
its \(s|U|\) units are in the neighboring columns, whose total over all
rows is \(s|N(U)|\). Hence \(|U|\le|N(U)|\), and the constructive
matching theorem supplies a permutation \(\pi\) supported on positive
entries. Subtract from each selected entry the minimum \(w\) of the
\(n\) selected entries. Every row and column sum decreases by \(w>0\);
at least one positive cell disappears, and none is added. Record
\((w,\pi)\). After at most \(n^2\) productive steps the residual is zero.
The recorded weights sum to \(D\) by any row sum, and undoing the
subtractions reconstructs \(DA\).

Each permutation matrix \(P_\pi\) uniquely maximizes
\(\sum_i A_{i,\pi(i)}\) with value \(n\): the entries are bounded by
one, and reaching \(n\) forces every selected entry to be one, exhausting
each row. Thus \(P_\pi\) is extreme. A nonpermutation matrix has a
decomposition involving two different permutations with positive weights.
Separating one of them from the normalized average of the rest gives
two distinct points: a positive-weight different permutation has a
positive coordinate outside the first one's support. This is a
nontrivial convex decomposition. Finally any linear objective on \(A\)
is the corresponding weighted average of its permutation values, and
an optimum permutation is obtained by a finite maximum.
\end{proof}

\begin{proposition}[Dimension, facets, and edges of the Birkhoff polytope]\label{result55}
The dimension is \((n-1)^2\). For \(n\ge3\) there are \(n^2\) facets,
one for every vanishing entry. Two different vertices form an edge
exactly when their permutations differ by a single nontrivial cycle.
Consequently the graph of the \(3\times3\) polytope is complete on six
vertices. The general number of edges is
\[
 \frac{n!}{2}\sum_{j=2}^n\binom nj(j-1)!.
\]
\end{proposition}
\begin{proof}
Every difference matrix has zero row and column sums, so its upper
left \((n-1)\times(n-1)\) square determines it. Given an arbitrary
rational square of that size, append negatives of its row and column
sums, then put the sum of the square in the bottom right. This
extends it uniquely to a zero-sum difference. Each coordinate unit
in the square thus yields a four-entry difference \(E\). At the uniform
matrix \(A_{ij}=1/n\), adding \(E/(2n)\) preserves nonnegativity.
These independent perturbations prove the dimension claim.

For a specified zero entry, relabel it as the bottom right. There is a
point on its hyperplane with all other entries positive when \(n\ge3\):
set the remaining entries of its row and column to \(1/(n-1)\) and the
upper square to \((n-2)/(n-1)^2\). For differences in this hyperplane
the sum of the upper square is zero; all but one of its entries are
free, and the remaining entry is minus their sum. Each of the resulting
basis differences, scaled by a rational smaller than the least positive
matrix entry divided by one plus its largest absolute coordinate, gives
a valid perturbation. The dimension drops by one. Every proper exposed
face is contained in some such hyperplane. Indeed take the average of
its maximizing permutations. If all entries of the average were positive,
each of the independent perturbations above from the full space could
be made in both directions by choosing a sufficiently small rational
step. Maximality would make the objective constant in all those directions
and hence constant on the polytope. A zero entry of the average is zero
in every maximizing permutation. The exhibited positive points outside
the single specified zero distinguish the \(n^2\) hyperplanes and prevent
one of their faces from containing another. Thus these are exactly the
maximal proper exposed faces.

For two different permutations, relabel columns to make the first the
identity. On each nontrivial cycle of the second permutation, a matching
inside the union of their supports has just two choices. Selecting one
off-diagonal edge occupies its successor column, forcing the off-diagonal
choice at the successor row and then around the cycle. Otherwise all
rows on the cycle use their diagonals. Fixed rows are forced, and
choices on different cycles are independent. With one cycle exactly
two vertices maximize the negative sum of entries outside the union.
They expose an edge. With two or more cycles, exchange one cycle between
the two original permutations to obtain two others \(R,S\) with
\(P+Q=R+S\). Any functional maximal at \(P,Q\) must be maximal at both
\(R,S\); neither lies on the segment from \(P\) to \(Q\), since an
interior point of that segment has fractional coordinates on the
nontrivial cycles. Hence the original pair is not an edge. There are
\(\binom nj(j-1)!\) ways to choose a \(j\)-cycle by selecting its support
and listing its cyclic order starting at its least label. Sum these
degrees and divide by two. Every nonidentity permutation of three
letters has precisely one nontrivial cycle.
\end{proof}

\subsection{The permutohedron and three elementary face formulas}

\begin{theorem}[Permutohedron]\label{result56}
For \(n\ge1\) the convex hull of the permutations of \((1,\ldots,n)\)
has a nonempty exposed face for each ordered partition of the coordinate
labels. A partition with \(k\) nonempty blocks gives a face of dimension
\(n-k\). If \(S(n,k)\) is the number of partitions of \([n]\) into
\(k\) unordered, nonempty blocks, the number of \(j\)-faces is
\((n-j)! S(n,n-j)\), \(0\le j<n\). Its \(n!\) vertices have degree
\(n-1\), and for \(n\ge2\) there are \(2^n-2\) facets.
\end{theorem}
\begin{proof}
For an objective with coefficients \(c_i\), group equal coefficients
and put the blocks in increasing order. In a maximizing permutation,
\(c_i<c_j\) forces \(x_i<x_j\), since swapping a misplaced pair
increases the objective by \((c_j-c_i)(x_i-x_j)>0\). At least one
permutation maximizes it by finite search. Every permutation obeying
the constraint assigns the same consecutive interval of values to
each block, and hence has the same objective as a maximizer. Thus these
are all the maximizers. Conversely, prescribe block indices as integer
coefficients to realize any ordered partition. The block order is
recovered from its maximizers: in different blocks the coordinate
comparison is constant across them, and within one block both orders
occur for any two distinct coordinates, by swapping their values.
This also gives the inclusion rule: splitting blocks into consecutive
subblocks restricts a face, and every nested pair is obtained this way.

The sum of the coordinates in each block is fixed. Thus one coordinate
per block is determined by the others. Fix a maximizing permutation and
an anchor label in each block. Swapping the anchor with each other label
in turn gives nonzero multiples of the independent coordinate differences
\(e_i-e_{\rm anchor}\), spanning the space of blockwise zero-sum
differences. This proves the dimension. Ordering the \(k\) blocks
multiplies \(S(n,k)\) by \(k!\). An edge has one block of size two and
all others singleton, so each vertex admits \(n-1\) swaps of consecutive
values. A facet has two ordered blocks, specified by a nonempty proper
first block.
\end{proof}

\begin{proposition}[Simplex, cube, and cross-polytope]\label{result57}
In dimension \(d\), the standard simplex has \(\binom{d+1}{j+1}\)
faces of dimension \(j\); the cube has \(2^{d-j}\binom dj\). For
\(d\ge1\) the cross-polytope has \(2^{j+1}\binom d{j+1}\) faces of
dimension \(j<d\) and one face of dimension \(d\). Each of these
polytopes and the permutohedron has
\(\sum_j(-1)^j f_j=1\), counting the whole polytope.
\end{proposition}
\begin{proof}
The vertices of the simplex are \(e_1,\ldots,e_{d+1}\). An objective's
maximum selects exactly the indices with largest coefficient. Any
nonempty set of indices is obtained by giving those indices coefficient
\(1\) and the others zero; a support of \(j+1\) vertices has dimension
\(j\) by subtracting a fixed vertex. In the cube with vertices
\(\{0,1\}^d\), each coordinate is fixed at \(0\) or \(1\) according
to the sign of its nonzero coefficient, or is free if that coefficient
vanishes. Every such pattern occurs; the free coordinate differences give
its dimension. The hull indeed contains every rational point with
\(0\le x_i\le1\): give vertex \(v\) the product of the \(x_i\) for
its ones and \(1-x_i\) for its zeros. These weights sum to one and have
barycenter \(x\) by distributivity.

The cross-polytope has vertices \(\{\pm e_i\}_{i=1}^d\). For a nonzero
objective take the largest absolute coefficient. The maximizing vertices
are the labels attaining it, each with its prescribed sign, so the
proper faces have one sign on a nonempty support. Conversely assign
these signs to its support and zero elsewhere. Any \(j+1\) selected
vertices are affinely independent by their distinct coordinate positions.
The zero objective yields the full \(d\)-dimensional hull, since differences
of opposite vertices span the whole coordinate space. These descriptions
give the counts.

For the simplex the alternating sum follows from \((1-1)^{d+1}=0\);
for the cube it is \((2-1)^d=1\). For the cross-polytope its sum over
proper faces is \(1-(1-2)^d\); add \((-1)^d\). For the permutohedron
put \(a_n=\sum_k(-1)^k k!S(n,k)\). Inserting a new label into one of
\(k\) blocks or making it a singleton gives
\(S(n+1,k)=kS(n,k)+S(n,k-1)\). The two sums contributing to \(a_{n+1}\)
combine coefficientwise as \(k-(k+1)=-1\), so \(a_{n+1}=-a_n\).
As \(a_0=1\), its face alternating sum \((-1)^na_n\) is one.
\end{proof}

\subsection{Kissing numbers without trigonometry}

A finite kissing configuration in dimension \(d\) is a list of vectors
whose squared norms have a common positive value \(L\) and whose pairwise
squared distances are at least \(L\). In ordinary Euclidean geometry these
are the centers of equal spheres of diameter \(\sqrt L\) touching the
one centered at the origin. The arithmetic predicate uses only products,
sums and comparisons.

\begin{theorem}[Small dimensions and explicit high-dimensional witnesses]\label{result58}
Over the ordered quadratic field \(K=\mathbb Q(\sqrt3)\) the maximum
sizes in dimensions one and two are \(2\) and \(6\), respectively.
Integer-coordinate configurations give lower bounds \(2d\) in every
positive dimension and \(2d(d-1)\) for \(d\ge2\). There is a rational
configuration of size \(240\) in dimension eight.
\end{theorem}
\begin{proof}
Write \(\theta^2=3,\ \theta>0\). For coding \(K\), use pairs of reduced
rationals \(a+b\theta\), with multiplication reduced using \(\theta^2=3\).
The norm \(a^2-3b^2\) vanishes only at zero: reducing a hypothetical
fractional square root \(p/q\) to lowest terms would give \(3\mid p^2\),
hence \(3\mid p\) by residues; substitution would also give \(3\mid q\).
Conjugation divided by the norm supplies an inverse. The positive cone
consists of quotients of two nonzero expressions with nonnegative
coefficients. This cone is closed under sums and products and disjoint
from its negative, by cross multiplication of expressions with nonnegative
coefficients. It determines a decidable sign: coefficients of like sign
are immediate; for opposite signs multiply by the positive
\(|a|+|b|\theta\), reducing the result to the rational \(a^2-3b^2\)
or its negative. Thus all comparisons below have primitive recursive
tests.

In dimension one \(u^2=v^2\) gives \((u-v)(u+v)=0\); at most two
different coordinates occur, and \(1,-1\) attain the bound. In dimension
two use the norm-preserving matrix
\[
 R=\frac12\begin{pmatrix}1&-\theta\\ \theta&1\end{pmatrix},
 \qquad R^3=-I,\qquad R^6=I .
\]
The six powers of \(R\) applied to \((1,0)\) have norm one and their
pairwise squared distances are \(1,3,4\), according to the distance
between their exponents modulo six. This gives six neighbors. To bound
an arbitrary list, divide directions into the six half-open cones obtained
by applying these powers to \(x>0,\ 0\le y<\theta x\).
Here is a signed comparison proof that the cones partition the nonzero
vectors. In the half-plane \(y>0\) together with the positive horizontal
axis, the inequalities \(\theta x>y\), \(-y<\theta x\le y\), and
\(\theta x\le-y\) are the three cases. Applying respectively \(I\),
\(R^{-1}\), and \(R^{-2}\) sends each case into the first cone,
by multiplication with the displayed matrix; the reverse substitutions
give precisely the same three cases. Negating handles the complementary
half-plane. This settles all boundary rays.

Two points in the same cone can be rotated into the first, preserving
norm and dot product. Write their slopes \(0\le a\le b<\theta\); their
coordinates have forms \((x,ax),(u,bu)\), with \(x,u>0\).
The elementary identity
\[
 4(1+ab)^2-(1+a^2)(1+b^2)
   =(3-b^2)+a(8b-a)+3a^2b^2>0
\]
uses only \(0\le a\le b\) and \(b^2<3\).
Multiply by \(x^2u^2\). The two points have common squared norm \(L\),
so their positive dot product exceeds \(L/2\). Their squared distance
is therefore less than \(L\). Each cone contains at most one center;
the six-column count gives the upper bound. The same algebra proves the
usual Euclidean inequality for a supplied list of real vectors; the
internal statement here quantifies over the codes in \(K\).

The \(2d\) signed unit vectors have squared norm one and pairwise
squared distances at least two. For the second family, take all
\(\pm e_i\pm e_j,\ i<j\), of norm squared two. Distinct such vectors
have inner product at most one: different supports meet in at most one
position, and a sign difference on an equal support brings the product
down to zero or below. Hence their squared distances are at least two.
This gives, among other cases, twelve centers in dimension three and
twenty-four in dimension four. In eight dimensions the second family
has \(112\) vectors. Add all \(\frac12(\pm1,\ldots,\pm1)\) having an
even number of minus signs. The first seven signs freely determine the
eighth, giving \(128\) new vectors of squared norm two. They have
inner product at most one with every two-coordinate vector. Two
different new sign patterns differ in a positive even number \(h\ge2\)
of positions and have inner product \(2-h/2\le1\). The \(240\) vectors
therefore meet the separation requirement.
\end{proof}

\paragraph{Formalization of the geometric statements.}
Use \texttt{matrix n D a} for a base-\((D+1)\) table of \(n^2\)
nonnegative numerators; the validity test checks each row and column
sum against \(D\). The Birkhoff loop stores a residual and a list of
integer weights and permutations. For at most \(n^2\) iterations maintain
the uniform residual sum, coordinatewise reconstruction of the input,
and the upper bound on remaining productive iterations from the support
size. Implement \texttt{cycleSupport} by following a permutation for
at most \(n\) steps. The test for the edge characterization counts its
nontrivial cycles; the proof constructs the two recombined vertices
from a cycle mask. For a coordinate objective define
\texttt{equalWeightBlocks} by sorting indices and grouping ties;
\texttt{maximizerMask} is a scan of all permutation codes. Prove the
strict swap inequality, its conclusion for any maximizing permutation,
and the equal objective of all block-respecting permutations. Encode
faces by the resulting masks; dimensions follow from the explicitly
listed independent differences and their coordinatewise spanning maps.
The same definitions work on the cube's \(2^d\) vertices or the other
finite lists. For kissing configurations define \texttt{Quad} using four
integer numerator/denominator fields (normalize denominator signs), and
prove the arithmetic and order rules from the conjugation and cone
calculations above. Implement \texttt{sector} by the three comparisons
and a sign bit. The squared-distance test and membership in a sector
are arithmetic comparisons with free coded inputs. Apply the finite
pigeonhole lemma to the resulting labels in \([6]\). A bounded
checker on all index pairs independently verifies each explicit lower
bound construction.
